\cotefile{15}
\scanpage{195}
2°) $X_{\deleted{\unsure{\text{réd}|\text{red}}}}$ domine $\operatorname{Spec}\mathbb{Z}$, i.e. $X$ contient des pts de car résiduelle nulle.

Alors \deleted{en plus de ces \inserted{dispose \deleted{\unsure{des}}} (et essentiellement \unsure{(i)s|(i)} (iii) dans le cas d’inégales car) on a, pour tt pt} \deleted{\unsure{(iii)} (iii) bis)}

(i) $\unsure{\xi|\zeta}_{\ell}\in \unsure{\mathbb{T}|T|\mathbf{T}}_{\unsure{\alpha|x}}(\mathbb{Z}_{\ell})$ \uline{tt} $\ell\in\mathbb{P}$, $\xi_\ell$ provenant sans doute de $T_\alpha(\widetilde{\mathbb{Z}}_{\ell\mathbb{Z}})$

$\xi_\ell$ dépend du choix d’un pt géom quelconque de $X$, et est indép. de ce choix modulo opération de $\unsure{\Pi|\pi}_1$, notation $\xi_{\bar{x},\ell}$ ou $\xi_{x,\ell}$

(ii) Si $X$ a des points $x$ de car. $p>0$, alors on dispose également des $\xi_{x,p}$, $\xi_{x,p}^{0}$.

\unsure{Correction|Connexion|Connection|Conjecture} : relation, \deleted{\unsure{de la}} si $y$ est un pt de car. $p>0$ et \deleted{\unsure{x}} un pt de car. $0$, entre $\xi_{y,p}$ et $\xi_{x,p}\otimes_{\mathbb{Z}_p}\mathrm{W}(k)$ [$y$ pt à val. ds corps parfait $k$]

\uline{Ils ne sont en général pas isomorphes}, sans doute. \deleted{\unsure{Ils}} Si \deleted{\unsure{ne sont pas}} $k$ est fini, ils sont cependant \inserted{isomorphes \unsure{par Lang|par Lang,}} \deleted{\unsure{et les groupes devant}} \deleted{\unsure{groupes de ces objets}} \deleted{\unsure{devant}} \deleted{\illegible{1}}

(iii) Si $x$ est un pt à val. dans un corps \uline{de car $0$}, on a $\xi_{x,\infty}\in T_\alpha(k)$ par Hodge-de Rham. \inserted{Avec ses \textcircled{\uuline{2}} filtrations \deleted{\unsure{et le groupe}}} Il n’y a pas lieu de

% margin note written sideways
\margin{L’une est celle qui donne comme gradué associé les $\unsure{H^{q}|H^{i}|H^{\bullet}}(x,\Omega^{\unsure{p|j|P}})$, l’autre, (\unsure{arithmétique|authentique|carthentique}) fait pendant à la filtration canonique des motifs. (Elle est \uline{triviale} pour les motifs simples.) La première filtration est originale au foncteur Hodge-De Rham, bien faire attention qu’elle ne provient pas d’une filtration du motif, c’est une filtration du \uline{foncteur}\ldots. !}

% archive sheet number: 194
